\chapter{Conclusions}
\label{ch:conclusions}

\section{What the thesis found, question by question}
\label{sec:concl:summary}

This thesis treated a decentralised prediction market as a data-mining object and asked four
questions (\cref{sec:intro:questions}): what its price process looks like and how its
trading is organised, which public information predicts its direction, which families of
models can read it, and how the answers change with the horizon.  The answers, on three
weeks of complete order-book and on-chain data, are these.

\paragraph{The process (\ref{q:process}).}  The price is a bounded, discrete process whose most
regular feature is its own activity clock: whether it will move is written in how recently and
how often it has moved, and the direction, when it moves, is weakly in the first level of the
book and in the reversal of its recent change.  The price level and the kind of contract
condition everything; the neighbours of a market, the wallets that trade it and the contracts
of its event, carry nothing that its own book does not already show.  Read as a graph, the
trading is one connected component of 646 thousand wallets and 158 thousand markets with a
dense core of a few thousand nodes, a long tail of wallets with a single fill, and a map
organised by category, contract type and date.  The sixty-one features
that summarise these properties are about seventeen independent quantities; the space they span
is organised by market and by activity, not by the label; the part of the label that is
linear or local is \emph{whether} the price moves, and the information the label asks for
moves from the activity clock at minutes to the shape of the book at hours.

\paragraph{The information and the models (\ref{q:information}, \ref{q:models}).}  A
stateless learner on sixty-one hand-made features that encode these properties is the strongest model at every horizon
from fifteen minutes to twelve hours, six to seven points above the best rule; a properly
trained tabular network ties it; a continuous-time stateful model reaches it only when given
the same inputs and the trees as a base, and beats it, by a small margin with an interval that
excludes zero, only at one minute, where the sequence of book events says something about the
next minute of a quiet market that no summary of the last hour contains.

\paragraph{The horizon (\ref{q:horizon}).}  At one minute the task is activity and the
sequence of book events carries information that the summary loses; at fifteen minutes and
one hour the summary is complete and a stateless learner is the right model; at six and
twelve hours the task is the direction of a slow reversal read with the price level and the
lifecycle of the market, and the stateless learner is six points above the rule while the
stateful models of this study are below it.  The live-game regime and the crypto interval
contracts are the hard cases at every horizon.

\paragraph{In more detail.}  The thesis built, from raw order-book and on-chain data, a leakage-free benchmark for
short-horizon price-direction prediction on Polymarket, and used it to compare a
family of models against the strongest simple rule on the data.  A continuous-time
model that contains the mean-reversion rule as a special case, encodes the order book
with a recurrent network, keeps a three-rate hierarchical memory and reads five
lookback price changes beats the rule by $2.4$ balanced-accuracy points at one hour and
$1.5$ at six hours, with a gain that is stable across seeds at one hour; reading the
order-book window one quantity at a time before mixing quantities (the time-first
encoder) raises the one-hour gain to $3.3$ points, reproduced on a second seed, and
the six-hour gain to $3.3$ points; at twelve hours the model is below the rule on the
full stream and above it on independent samples.  The
wallet--market transaction graph never adds a
reliable gain; zero-shot time-series foundation models re-discover the rule and do not
beat it, and as inputs of the best model they add between three and six tenths of a
point at one hour over two seeds, a gain confined to central prices.  The last experiment
turns the picture around: a gradient-boosted model with no state, on sixty-one hand-made
features from the clean data, beats the best deep model by $3.7$ points at one hour,
$3.2$ at six hours and $8.0$ at twelve, and is $7.3$ points above the rule at one minute;
the features that carry the gain (the time since the last price change, the number of
recent changes, the first level of the book, the long returns read together with the
price level) are exactly the ones the event-clock tensor hides or the linear direct path
cannot use.  Given those same features as inputs, the deep model recovers half of the gap
at one hour ($0.584$); the other half is the cost of its online, time-ordered training regime,
and with the cross-fitted trees as a base and the same calibration as the trees it ties them.
At one minute, with the control rebuilt on the runs' own samples and every model calibrated
the same way, the stateful model is $+0.6$ above the trees alone and $+1.4$ as their residual;
at fifteen minutes it ties them.  The locked test days of 24--26 March reproduce every one
of these orderings, with the one-minute margins at $+0.4$ to $+0.8$ alone and $+0.9$ to $+1.0$ as the
residual, over two seeds each, and the long-horizon
gaps one to two points wider (\cref{sec:horizons:test}).

\section{Limitations}
\label{sec:concl:limits}

\begin{itemize}[nosep]
  \item One validation day and one training week.  The three locked test days of 24--26
  March (\cref{sec:horizons:test}) reproduce every ordering of the validation day, but they
  are three days of the same month, and the long-horizon gaps widened by one to two points
  on them.
  \item Balanced accuracy is a statistical, not an economic, measure: a one-tick move
  predicted correctly is not a profitable trade once spread and fees are paid.  The
  thesis makes no claim about trading.
  \item The graph was tested with one family of message-passing operators and a shared
  decay; a wallet memory at a slower time scale remains untested.
  \item Hyper-parameters were chosen by a small number of screening runs, not by a
  systematic search; the comparison between components is paired and fair, the absolute
  level could be higher.
  \item Three weeks of one venue in one month: the regularities found are those of
  Polymarket in March 2026, and a different season of events, or a different exchange,
  could weigh the families differently.
\end{itemize}

\section{Towards deployment}
\label{sec:concl:deployment}

The thesis stops before deployment.  Balanced accuracy is a statistical measure; a trading use of these models would
have to answer at least four further questions, none of which the present data settle.
\emph{Costs}: most correctly predicted moves are one tick, and a tick is smaller than the
median spread of three of the four families of markets of \cref{tab:structure:clusters}, so a correct direction is not a
profitable trade unless the position is taken as a resting order, which changes the
problem.  \emph{Latency}: a prediction at block $t$ uses the book up to $t$ and the fills of
$t$; the order that acts on it reaches the operator later, and at one minute the signal is
mostly ``the price is about to move'', which a resting order cannot use.  \emph{Capacity}:
the sparse markets, where the sequential state is worth most, are sparse because little
money trades them.  \emph{Non-stationarity}: the families were selected on one validation
day; a deployed model would be retrained on a rolling window, and the stability of the
rankings across days is measured here only by the locked test period.  What the thesis does
settle is what a model of this market should be made of, and that is the content of the next
section.

\section{What the thesis would start from today}
\label{sec:concl:today}

Taken together, the control suggests the design the thesis would start from today: a
regular time grid per market (one-minute bars, aggregated to fifteen minutes for the long
horizons) instead of the event clock, a feature set reduced to the first level, the
microprice, the recent-mean distance, the activity clock and the contract type, the
time-first encoder on those bars, and gradient boosting as the control that any deep model
must beat.

\section{Future work}
\label{sec:concl:future}

\begin{itemize}[nosep]
  \item A rebuilt input pipeline on a regular time grid (one-minute bars, fifteen-minute
  bars for the long horizons) with the features the no-state control found decisive, the
  time-first encoder on those bars, a learned contract-type embedding and one
  multi-horizon model \cite{zhang2021multihorizon}; gradient boosting as the control
  that every deep model has to beat; and the residual construction at fifteen minutes and
  at the long horizons, where it was not run.
  \item A learned, sparse market--market graph restricted to markets of the same event,
  in the spirit of \textcite{cini2023sparse}, in place of the wallet graph.
  \item Smoothed labels at long horizons \cite{berti2025tlob}, to make the
  \textsc{const} class less sensitive to one-tick noise.
  \item The Binance price as a covariate for the crypto markets, following
  \textcite{young2026openmarket}.
  \item Foundation-model features (\cref{sec:models:tsfm}) at the horizons where the
  model is already strong, and a fine-tuned or domain-pretrained variant if the frozen
  one helps.
  \item The time-first encoder with the elapsed time between snapshots as an input
  channel, and the time-first encoder together with the foundation-model features,
  since the two gains are measured on the same base model and may add up.
  \item A different task on the same data: the generation of the next edges of the
  transaction graph (which wallet will trade which market), the estimation of the real
  probability of the event against the price, or the distribution of the price at the
  horizon rather than its sign; and a graph model with a wallet memory at the time scale of
  days, which \cref{sec:expl:graph} suggests is the only scale at which most wallets have a
  history.
  \item An ensemble of the seeds already trained, the cheapest known improvement of a
  deep classifier \cite{lakshminarayanan2017ensembles}, and event-based labels such
  as the triple-barrier method \cite{lopezdeprado2018advances} in place of the
  fixed-horizon label.
\end{itemize}

